\documentclass[10pt]{article}

\usepackage{tmlr}
\usepackage{amsmath,amssymb}
\usepackage{booktabs}
\usepackage{microtype}
\usepackage{tabularx}
\usepackage{xcolor}
\usepackage{hyperref}
\usepackage{url}

\hypersetup{colorlinks=true,citecolor=blue,linkcolor=blue,urlcolor=blue}
\newcommand{\arca}{\textsc{ARCA}}
\newcommand{\system}[1]{\textsf{#1}}
\newcommand{\code}[1]{\texttt{#1}}

\title{When Successful Runs Do Not Mean Valid Learning:\\
Auditing Silent Reward Contracts in Agentic Reinforcement Learning}
\date{}

\begin{document}
\maketitle

\begin{abstract}
Reliable agentic reinforcement learning (RL) depends on more than a job
exiting successfully.  Its pipelines connect multi-turn traces, tool
results, evaluators, rewards, group advantages, token masks, and experiment
artifacts.  A job can finish normally while one of these boundaries silently
changes the treatment or erases the intended signal.  We call this a
\emph{silent reward-contract failure} and present \arca, a CPU-only preflight
auditor for typed trace equivalence, action--result pairing, reward flow and
informativeness, grouped-advantage support, mask identity, status, and
provenance.  The study separates discovery from frozen evaluation.  Six
natural cases are reproduced across five source-pinned systems; two change a
registered scientific conclusion.  On an unseen OpenRLHF revision, frozen
rules detect all 40 registered single-fault mutations across four families,
while flagging none of 40 source-derived clean controls; native preflight
detects 25\% of the mutations.  A later five-system purposive validation finds
two additional natural cases and yields 7/7 exact case predictions with 0/5
clean false positives, whose Wilson 95\% interval is nevertheless
$[0.0000,0.4345]$.  We therefore claim case-based external transfer, not
ecosystem prevalence or reliability certification.  A preregistered GPU
attempt failed before model loading and is reported as infrastructure-invalid,
not as evidence about learning.  The contribution is an executable audit
methodology, a source-pinned case corpus, and fail-closed reporting practices
for establishing whether an agentic RL experiment implemented its intended
reward contract.  Our results motivate contract preflight as a validity check
before learning curves are interpreted, while leaving downstream model effects
for a separately valid experiment.
\end{abstract}

\section{Introduction}

Multi-turn agent training is no longer a single function from prompt to scalar
reward.  A rollout may cross a serializer, tool protocol, environment adapter,
reward worker, grouping rule, token selector, and asynchronous storage layer
before it becomes a gradient update.  Recent work improves reward design
~\citep{zhang2026cm2,modecrua2026irc}, credit assignment
~\citep{zhang2026credit}, and training stability~\citep{wang2025ragen}.  These
advances generally assume that the value called ``reward,'' the trajectory
being scored, and the tokens being optimized still represent the intended
experiment.

That assumption can fail without a crash.  A JSON object can become a quoted
string, a missing reward can become the valid scalar zero, a producer can write
one reward field while a consumer reads another, or two nominal mask arms can
select the same response tokens.  Standard checks---successful exit, finite
loss, schema validity, or non-empty output---need not expose these errors.  The
run is operationally successful but scientifically uninterpretable.

We study this earlier validity question:

\begin{quote}
Did the executed pipeline preserve the intended trace, reward, comparison,
treatment, and evidence contracts strongly enough for the resulting learning
claim to be meaningful?
\end{quote}

We make four contributions.  First, we define a contract chain and eight
executable failure rules spanning six families.  Second, we provide \arca, a
small framework-independent auditor with explicit positive controls.  Third,
we evaluate frozen rules on source-pinned natural cases, registered mutations,
an unseen framework, and a later purposive five-system sample.  Fourth, we
release a fail-closed artifact that preserves failed and unopened runs instead
of selectively repairing or excluding them.

Our claim is deliberately narrower than a prevalence study or a new RL
algorithm.  The evidence shows that silent contract failures occur in multiple
public systems and that the registered checks transfer to specific unseen
boundaries.  It does not estimate how often such failures occur in the field,
show that audited projects are broadly unreliable, or demonstrate improved
model quality.

\section{Related Work and the Remaining Gap}

\paragraph{Reward quality and reward exploitation.}
Checklist rewards and iterative calibration make feedback more informative
for multi-turn tool use~\citep{zhang2026cm2,modecrua2026irc}.  Reward-hacking
benchmarks instead test whether agents exploit evaluators or task shortcuts
~\citep{thaman2026rhb}.  Both ask whether the objective is useful or gameable.
Our failure unit is different: an honest trajectory and intended evaluator can
be corrupted by a software boundary before optimization.

\paragraph{Credit assignment and agentic RL systems.}
Credit-assignment work distributes outcome information across tokens, steps,
turns, or agents~\citep{zhang2026credit}.  RAGEN and Agent-R1 expose
trajectory and optimization abstractions~\citep{wang2025ragen,cheng2025agentr1}.
Agent Lightning and OpenRLHF separate execution, reward, and training
interfaces~\citep{luo2025lightning,openrlhf2026}.  veRL and AReaL provide
multi-turn or asynchronous training boundaries~\citep{verl2026,areal2026};
slime and rLLM provide further reward and evaluator interfaces
~\citep{slime2026,rllm2026}.  Modularity creates explicit boundaries but does
not itself test cross-boundary semantic preservation.

\paragraph{Pipeline evaluation.}
Agent$^2$ RL-Bench evaluates whether coding agents can build and close an RL
loop under a budget~\citep{chen2026agent2}.  \arca{} instead evaluates a
pipeline supplied by a researcher: its unit is a typed contract and a positive
control, not the engineering performance of an autonomous coding agent.

The closest methodological comparators are source-framework preflights and
Agent$^2$ RL-Bench.  Native preflights validate one framework's accepted
inputs, while Agent$^2$ RL-Bench scores whether an engineering agent closes an
RL loop.  Neither evaluates whether a supplied experiment preserves typed
trace meaning, reward flow, advantage support, and treatment identity across
framework boundaries.  Conversely, \arca{} does not evaluate autonomous
engineering ability or propose a training algorithm.  This distinction is the
claimed novelty delta.

Our refreshed search over arXiv, OpenReview, official documentation, and
official repositories applied five substitute criteria: typed trace checks,
reward and advantage checks, treatment-mask identity, cross-framework
execution, and held-out framework evaluation.  No screened work satisfied all
five.  The remaining gap is therefore a frozen, executable preflight for
semantic experiment contracts, evaluated on source-pinned cases and an unseen
framework.  This is a bounded search conclusion, not proof of universal
novelty.

\section{Reward Contracts}

Let an intended experiment be the contract chain in
Equation~\ref{eq:contract_chain}:
\begin{equation}
  x \xrightarrow{S} \tau \xrightarrow{E} r
  \xrightarrow{G} A \xrightarrow{M} z
  \xrightarrow{L} \mathcal{A},
  \label{eq:contract_chain}
\end{equation}
where $S$ serializes an interaction trace $\tau$, $E$ evaluates it, $G$ forms a
comparison or advantage $A$, $M$ selects trainable tokens $z$, and $L$ writes
the evidence artifact $\mathcal{A}$.  A contract $C_j$ specifies an invariant
at one boundary and an executable positive control that must distinguish at
least two intended inputs.

We define a \emph{silent reward-contract failure} as a violation of some
$C_j$ for which the enclosing process can still terminate successfully and
emit superficially valid outputs.  This excludes loud exceptions that remain
visible in structured status.  It also excludes explicit, documented fallback
behavior when its occurrence remains observable to downstream analysis.

Table~\ref{tab:taxonomy} lists the registered families.  In the stable codes,
we use typed-argument equivalence (TYPE), reward transport (REWARD), flow
continuity (FLOW), and not scientifically evaluated (NOT).  The implementation returns
these codes rather than framework-specific error messages.

\begin{figure*}[t]
\centering
\fbox{\parbox{0.94\textwidth}{\centering
Prompt and tool state $\rightarrow$ typed trace $\rightarrow$ evaluator reward
$\rightarrow$ grouped advantage $\rightarrow$ trainable-token mask
$\rightarrow$ evidence artifact\\[3pt]
\small F1 checks trace meaning; F2 checks pairing and reward flow; F3--F4
check signal and comparison support; F5 checks treatment identity; F6 checks
status and provenance.}}
\caption{The audited contract chain from executed interaction to scientific
artifact.  A process-level success signal can coexist with a semantic failure
at any boundary.}
\label{fig:contract_chain}
\end{figure*}

\begin{table*}[t]
\centering
\small
\caption{Registered contract families.  F2 and F6 each contain two executable
rules, giving eight rule codes across six families.}
\label{tab:taxonomy}
\begin{tabularx}{\textwidth}{@{}lllX@{}}
\toprule
Family & Rule & Boundary & Positive-control invariant \\
\midrule
F1 & TYPE & trace serialization & Equivalent tool arguments canonicalize to one JSON object. \\
F2 & PAIRING & trace assembly & Tool calls and results have equal identifier multisets. \\
F2 & REWARD-FLOW & reward adapter & Intended keys are both produced and consumed. \\
F3 & INFORMATIVE & evaluator output & Registered contrasting controls yield finite, distinct rewards. \\
F4 & ADVANTAGE & group construction & Every comparison group has at least two items and some reward support. \\
F5 & TREATMENT & token masking & Distinct, non-aliased arms do not select identical token masks. \\
F6 & STATUS & failure coercion & A numeric fallback retains a structured failure or missingness status. \\
F6 & PROVENANCE & artifact boundary & Observed source and protocol identifiers equal required identifiers. \\
\bottomrule
\end{tabularx}
\end{table*}

\section{ARCA Auditor}

\arca{} consumes one schema-valid case containing a framework, workload,
seed, evidence class, expected rule codes, and a payload.  Each rule is a
deterministic predicate.  For example, TYPE parses string arguments and
canonicalizes JSON objects; PAIRING compares identifier multisets;
REWARD-FLOW intersects intended, produced, and consumed keys; INFORMATIVE
requires finite contrasting rewards; and TREATMENT compares masks after
declared aliases are removed.  STATUS fires only when a failure is represented
as numeric reward without structured status.  PROVENANCE compares required
and observed identifiers.

The design is intentionally small.  It neither imports an RL framework nor
attempts to infer a user's objective from logs.  Framework adapters extract
facts from source-pinned boundaries into the common schema.  This separates
the auditable rule from the code needed to observe a particular system.

\paragraph{Fail-closed execution.}
Every generator records a protocol identifier, deterministic case identity,
source hashes, execution state, and finite measurements.  Duplicate steps,
unaligned arms, source drift, or frozen-evaluator drift invalidate an artifact.
Failures and unopened runs remain outcomes.  A controller may stop a bounded
experiment, but it may not silently repair and splice results into the consumed
protocol.

\section{Study Design}

The study uses four evidence classes and never pools their denominators.  These
classes instantiate the boundaries in Figure~\ref{fig:contract_chain}:

\begin{enumerate}
  \item \textbf{Natural cases}: behavior reproduced from a pinned production
  or documented execution boundary without inserting a synthetic fault.
  \item \textbf{Source-derived clean controls}: valid states constructed from
  the same inspected boundary.
  \item \textbf{Registered mutations}: one predeclared fault inserted into an
  otherwise clean case.
  \item \textbf{Post-freeze replications}: cases found after the primary rule
  and held-out evaluation were frozen.
\end{enumerate}

\subsection{Development and freeze}

The development set contains 163 cases.  Three are natural cases from AReno
and veRL; the remaining 160 are balanced clean controls and single-fault
mutations over eight rule codes and ten fixture seeds.  Before inspecting the
held-out framework, we froze the auditor, evaluator, and development cases by
the Secure Hash Algorithm (SHA)-256.  Development macro recall was 1.00.  The
clean false-positive rate (FPR) was 0/80 with Wilson
95\% interval $[0.0000,0.0458]$, and all three natural cases were detected.

\subsection{One-shot held-out evaluation}

The held-out system is OpenRLHF at commit
\code{bc71bb19464aca306b33080b2d2bb45d154e2f49}.  Its adapter was written only
after the freeze.  Inspection produced 40 source-derived clean controls and 40
registered single-fault mutations across reward flow, group advantage,
treatment identity, and failure status.  No natural OpenRLHF failure was found;
therefore this split measures mutation transfer, not natural-case prevalence.

We compare \arca{} with exit-only, finite-only, schema-only, and the native
preflight observable at the inspected OpenRLHF boundary.  The primary metrics
are macro recall across mutation families and clean false-positive rate.  We
report Wilson intervals and a preregistered hierarchical bootstrap over the
framework--workload--family cells; repeated fixture seeds are measurements,
not independent framework samples.

\subsection{External natural validation}

After the held-out evaluation, five previously unseen systems were selected
by frozen inclusion criteria before their relevant files were opened: slime,
Agent-R1, RAGEN, rLLM, and Agent Lightning.  All remained in the denominator
after inspection.  Qualifying natural behavior was dynamically reproduced by
compiling and executing only the pinned function bodies.  A framework-level
clean control was retained for every candidate, including negative
inspections.  This sample is purposive and small.

\section{Natural Case Studies}

Table~\ref{tab:natural} reports all six natural cases.  A conclusion flip means
that the observed behavior changes a preregistered treatment distinction or a
gate, not necessarily downstream model quality.

\begin{table*}[t]
\centering
\small
\caption{Natural cases, kept separate by boundary and evidence phase.}
\label{tab:natural}
\begin{tabularx}{\textwidth}{@{}llllXc@{}}
\toprule
Phase & System & Rule & Severity & Reproduced behavior & Flip \\
\midrule
Dev & AReno & F1 TYPE & high & Tool argument alternatives are not equivalent JSON objects. & yes \\
Dev & veRL & F2 REWARD-FLOW & high & Intended tool reward is produced under a field the consumer does not use. & yes \\
Dev & veRL & F6 STATUS & medium & Timeout is converted to numeric zero without structured status. & no \\
Post-freeze & AReaL & F6 STATUS & medium & Missing reward becomes zero in the assembled tensor without missingness status. & no \\
P5 & slime & F6 STATUS & medium & Missing evaluation-replay reward becomes 0.0 with no reward-present field. & no \\
P5 & rLLM & F3 INFORMATIVE & high & Dicts lacking \code{reward} map to 0.0 even when \code{is\_correct} differs. & yes \\
\bottomrule
\end{tabularx}
\end{table*}

The rLLM behavior occurs in two production coercion functions but is counted
once because both instantiate the same evaluator contract.  The slime case is
an official evaluation-replay boundary, not a core training-path finding.
Agent-R1, RAGEN, and Agent Lightning are bounded negative inspections: their
audited fallbacks retained structured errors, warnings, or reward-presence
evidence and were not relabelled as findings.

\section{Results}

\subsection{Frozen held-out transfer}

Table~\ref{tab:heldout} shows the one-shot result.  \arca{} detects 10/10
mutations in each of four held-out families, for macro recall 1.00.  It flags
0/40 source-derived clean controls; the Wilson 95\% interval for clean FPR is
$[0.0000,0.0876]$.  Exit-only, finite-only, and schema-only checks detect none
of the registered semantic faults.  The native preflight detects one family,
for macro recall 0.25.  The hierarchical bootstrap interval for \arca{} macro
recall is $[1.00,1.00]$ under the frozen registered cells; it does not cover
unregistered fault families.

\begin{table}[t]
\centering
\small
\caption{Held-out OpenRLHF mutation transfer.  Each method has clean FPR 0/40;
the shared point estimate must be read with its $[0.0000,0.0876]$ Wilson
interval.}
\label{tab:heldout}
\begin{tabular}{@{}lc@{}}
\toprule
Auditor & Mutation macro recall \\
\midrule
Exit-only & 0.00 \\
Finite-only & 0.00 \\
Schema-only & 0.00 \\
Native preflight & 0.25 \\
\arca{} & 1.00 \\
\bottomrule
\end{tabular}
\end{table}

\subsection{External validation}

Table~\ref{tab:external} retains every selected framework.  The five-system
external sample contains two natural cases and five clean
controls.  Frozen rules detect both natural cases, including the one registered
high-severity case, and produce the exact expected rule set on all 7 cases.
The clean FPR is 0/5.  Its Wilson 95\% interval is
$[0.0000,0.4345]$, which is too wide to support a claim that ecosystem FPR is
low.  The external result supports case-based transfer to two previously
unseen boundaries and documents three negative inspections; it is not a
reliability certificate.

\begin{table}[t]
\centering
\small
\caption{External validation retains all selected systems.}
\label{tab:external}
\begin{tabularx}{\columnwidth}{@{}lXl@{}}
\toprule
System & Audited boundary & Outcome \\
\midrule
slime & forge evaluation replay & natural F6 case \\
Agent-R1 & WebShop environment step & bounded negative \\
RAGEN & Lean timeout/server error & bounded negative \\
rLLM & Python evaluator coercion & natural F3 case \\
Agent Lightning & veRL daemon reward fill & bounded negative \\
\bottomrule
\end{tabularx}
\end{table}

\section{Invalid Dynamic Attempt}

A separately frozen GPU protocol was intended to test whether strict versus
canonical reward handling changes downstream training.  The first of six
commands failed after 12.8175 seconds because the reward module could not
import the project package.  The failure occurred before model loading,
rollout generation, reward evaluation, or any optimizer update.  Scientific
trajectories and training updates were both zero; the remaining five commands
were unopened.  Observed single-GPU occupancy was 0.00356 hours.

The protocol outcome and scientific status are, respectively,
\begin{quote}
\small\ttfamily
INVALID\_P4\_INFRASTRUCTURE\_REWARD\_IMPORT\\
NOT\_EVALUATED
\end{quote}
No P4 arm difference appears in our analysis.  This example also illustrates
why fail-closed artifacts matter: repairing and selectively rerunning the
consumed protocol would erase the actual execution history.

\section{Threats to Validity}

\paragraph{Selection and prevalence.}
Systems were purposively selected for relevant agentic-RL boundaries and
public inspectability.  Six cases across five systems establish existence and
heterogeneity, not prevalence.  The external 0/5 clean result has a 0.4345
upper confidence bound and cannot certify low FPR.

\paragraph{Mutation construct validity.}
The held-out score uses registered single-fault mutations.  These isolate rule
behavior but may not represent the joint faults, concurrency effects, or
version-specific interactions seen in production.  Perfect registered recall
is not universal coverage.

\paragraph{Source-to-runtime validity.}
Natural cases were tied to pinned source and, where feasible, exact function
bodies on CPU.  We did not run full distributed training for these inspections.
Compiler extraction can miss effects introduced by surrounding services.  We
therefore describe exact audited boundaries and avoid project-wide claims.

\paragraph{Rule dependence.}
Adapters encode which facts are observable.  A malformed adapter can hide or
create a violation.  Source hashes, common schemas, positive controls, and
held-out freeze reduce but do not eliminate this risk.

\paragraph{No downstream model claim.}
P4 generated no scientific trajectory.  We do not know the size or direction
of any model-quality effect.  The present contribution is preflight validity,
not improved optimization.

\section{Responsible Reporting and Broader Impact}

Publishing source-pinned software cases can help researchers avoid invalid RL
runs and wasted compute.  It can also be misread as a security disclosure or a
ranking of framework quality.  We mitigate this risk by reporting bounded
functions and commits, retaining negative inspections, separating explicit
fallbacks from silent failures, and avoiding vulnerability-rate language.
No private system, user data, model output, or credential is included.  Public
maintainer contact and disclosure are outside this artifact and require a
separate decision.

The auditor can itself create false confidence if users treat passing eight
rules as complete correctness.  The intended use is as an extensible preflight
with registered positive controls, followed by task-specific validation.  It
is not a general reward-hacking defense, formal verification system, or safety
certificate.

\section{Conclusion}

Agentic RL experiments can preserve operational success while losing the
scientific meaning of a trace, reward, comparison, or treatment.  \arca{}
turns a narrow set of these assumptions into eight executable rules.  Frozen
rules detected all 40 held-out mutations with 0/40 clean flags, and later
matched two natural cases plus five controls without rule changes.  These
results are bounded by mutation dependence, purposive system selection, a
wide external-control interval, and the absence of a valid downstream training
study.  Future work should test independently contributed contracts and, under
a separately frozen protocol, measure whether fail-closed preflight prevents
invalid runs at acceptable overhead.  Before interpreting a learning curve,
researchers should verify that contrasting traces reach contrasting rewards,
advantages, masks, and artifacts under pinned provenance.

\bibliographystyle{tmlr}
\bibliography{references}

\appendix
\section{Reproducibility Checklist}

The anonymous artifact contains the frozen evaluator, case generators, source
manifests, JSON and CSV results, exact CPU commands, and a deterministic archive
builder.  Every distributed file is allowlisted and hashed.  A verifier rejects
unexpected files, hash drift, non-parsing JSON, headerless CSV, local absolute
paths, identity strings, Secure Shell (SSH) endpoints, credential-like
assignments, and Git
metadata.  The invalid P4 evidence archive is deliberately excluded; only its
redacted gate summary is distributed.

\section{Interpretation Rules}

For a mutation family with $x$ detections out of $n$ registered cases we report
recall $x/n$.  For $k$ clean flags among $m$ controls we report $k/m$ and the
two-sided Wilson interval.  Framework--workload--family is the inferential
cell; fixture seeds are repeated measurements.  Missing, failed, or unopened
runs remain explicit terminal outcomes.  Natural and synthetic evidence is
never aggregated into a vulnerability rate.

\section{Source Pins}

The central frozen revisions are:
\begin{description}
  \item[veRL] \nolinkurl{e9618406de5bad40041d7612554e465ec2003ec1}
  \item[OpenRLHF] \nolinkurl{bc71bb19464aca306b33080b2d2bb45d154e2f49}
  \item[AReaL] \nolinkurl{62f955c5e0388aebc6fd58c5ad3f8fcf9d7384b8}
\end{description}
The external manifest records the pinned revisions and SHA-256 hashes for
slime, Agent-R1, RAGEN, rLLM, and Agent Lightning.  These identifiers define
the scope of every technical claim; later upstream behavior may differ.

\end{document}
